\textbf{Phần 1. Trắc nghiệm lựa chọn}
\setcounterpageref{ex}{0}

\begin{ex} %[2D12V62]
Kết quả của phép tính $632 - 275$ là
\choice
{$367$}
{$347$}
{$457$}
{\True $357$}
\loigiai{
$632 - 275 = 357$.
\textbf{Đáp án của câu hỏi là D}
}
\end{ex}

\begin{ex} %[2D12V62]
Tính hiệu của $800 - 345$.
\choice
{$465$}
{\True $455$}
{$555$}
{$445$}
\loigiai{
$800 - 345 = 455$.
\textbf{Đáp án của câu hỏi là B}
}
\end{ex}

\begin{ex} %[2D12V62]
Điền chữ số thích hợp vào dấu $*$ để: $7*4 - 289 = 445$.
\choice
{$2$}
{$4$}
{\True $3$}
{$5$}
\loigiai{
$445 + 289 = 734$, vậy $* = 3$.
\textbf{Đáp án của câu hỏi là C}
}
\end{ex}

\begin{ex} %[2D12V62]
Một kho thóc có $920$ bao thóc. Sau khi xuất đi một số bao thì trong kho còn lại $465$ bao. Hỏi kho đã xuất đi bao nhiêu bao thóc?
\choice
{\True $455$ bao}
{$445$ bao}
{$465$ bao}
{$555$ bao}
\loigiai{
$920 - 465 = 455$ (bao).
\textbf{Đáp án của câu hỏi là A}
}
\end{ex}

\begin{ex} %[2D12V62]
Tìm $x$ biết: $x + 286 = 614$.
\choice
{$318$}
{$338$}
{$428$}
{\True $328$}
\loigiai{
$x = 614 - 286 = 328$.
\textbf{Đáp án của câu hỏi là D}
}
\end{ex}

\begin{ex} %[2D12V62]
Một người có một mảnh đất dài $1 000\text{ m}$, người đó đã rào được $675\text{ m}$. Còn lại bao nhiêu mét đất chưa được rào?
\choice
{$315\text{ m}$}
{\True $325\text{ m}$}
{$335\text{ m}$}
{$425\text{ m}$}
\loigiai{
$1 000 - 675 = 325$ (m).
\textbf{Đáp án của câu hỏi là B}
}
\end{ex}

\begin{ex} %[2D12V62]
Một đàn vịt có $540$ con, trong đó có $285$ con vịt trắng, còn lại là vịt khoang. Đàn vịt có bao nhiêu con vịt khoang?
\choice
{$245$ con}
{$265$ con}
{\True $255$ con}
{$355$ con}
\loigiai{
$540 - 285 = 255$ (con vịt khoang).
\textbf{Đáp án của câu hỏi là C}
}
\end{ex}

\begin{ex} %[2D12V62]
Hiệu giữa số lớn nhất có ba chữ số và $678$ là
\choice
{\True $321$}
{$311$}
{$331$}
{$221$}
\loigiai{
$999 - 678 = 321$.
\textbf{Đáp án của câu hỏi là A}
}
\end{ex}

\begin{ex} %[2D12V62]
Điền dấu thích hợp: $524 - 186 \dots 612 - 274$.
\choice
{$>$}
{$<$}
{Không so sánh được}
{\True $=$}
\loigiai{
$524 - 186 = 338$; $612 - 274 = 338$. Hai vế bằng nhau.
\textbf{Đáp án của câu hỏi là D}
}
\end{ex}

\begin{ex} %[2D12V62]
Một cửa hàng có $700\text{ kg}$ gạo. Ngày đầu bán $265\text{ kg}$, ngày thứ hai bán $185\text{ kg}$. Cửa hàng còn lại bao nhiêu ki-lô-gam gạo?
\choice
{$240\text{ kg}$}
{\True $250\text{ kg}$}
{$260\text{ kg}$}
{$350\text{ kg}$}
\loigiai{
Đã bán: $265 + 185 = 450$ kg. Còn lại: $700 - 450 = 250$ kg.
\textbf{Đáp án của câu hỏi là B}
}
\end{ex}

\begin{ex} %[2D12V62]
Tìm số tự nhiên $x$ sao cho: $450 - x = 186 + 74$.
\choice
{$180$}
{$200$}
{\True $190$}
{$210$}
\loigiai{
$186 + 74 = 260$. $x = 450 - 260 = 190$.
\textbf{Đáp án của câu hỏi là C}
}
\end{ex}

\begin{ex} %[2D12V62]
Hiệu của hai số là $250$. Nếu bớt ở số bị trừ $30$ đơn vị và thêm vào số trừ $20$ đơn vị thì hiệu mới là bao nhiêu?
\choice
{\True $200$}
{$240$}
{$260$}
{$300$}
\loigiai{
Bớt số bị trừ 30 thì hiệu giảm 30; thêm số trừ 20 thì hiệu giảm tiếp 20. Hiệu mới: $250 - 30 - 20 = 200$.
\textbf{Đáp án của câu hỏi là A}
}
\end{ex}

\textbf{Phần 2. Trắc nghiệm đúng sai}
\setcounterpageref{ex}{0}

\begin{ex} %[2D12V62]
Xét tính đúng sai của các phép trừ có nhớ phạm vi 1 000:
\choiceTF
{\True 512 - 187 = 325.}
{700 - 250 = 550.}
{\True 1 000 - 555 = 445.}
{\True Khi thực hiện phép trừ có nhớ, ta mượn 1 ở hàng liền kề bên trái.}
\loigiai{
\begin{itemize}
    \item Ý a) Đúng: 12 - 7 = 5 nhớ 1; 11 - (8+1) = 2 nhớ 1; 5 - (1+1) = 3. Đúng.
    \item Ý b) Sai: 700 - 250 = 450, không phải 550.
    \item Ý c) Đúng: 1 000 - 555 = 445. Đúng.
    \item Ý d) Đúng: Quy tắc trừ có nhớ. Đúng.
\end{itemize}
\begin{itemize}
    \item a) \textbf{Đúng}.
    \item b) \textbf{Sai}.
    \item c) \textbf{Đúng}.
    \item d) \textbf{Đúng}.
\end{itemize}
}
\end{ex}

\begin{ex} %[2D12V62]
Một bể nước có 820 lít. Người ta dùng 350 lít để giặt và 280 lít để tưới rau:
\choiceTF
{\True Tổng số nước đã dùng là 630 lít.}
{\True Lượng nước còn lại trong bể là 190 lít.}
{Số nước giặt đồ ít hơn số nước tưới rau.}
{\True Nếu muốn bể đầy lại 820 lít thì cần bơm thêm 630 lít.}
\loigiai{
\begin{itemize}
    \item Ý a) Đúng: 350 + 280 = 630 l. Đúng.
    \item Ý b) Đúng: 820 - 630 = 190 l. Đúng.
    \item Ý c) Sai: 350 > 280 nên giặt đồ dùng nhiều hơn. Sai.
    \item Ý d) Đúng: Đúng vì đã dùng 630 lít.
\end{itemize}
\begin{itemize}
    \item a) \textbf{Đúng}.
    \item b) \textbf{Đúng}.
    \item c) \textbf{Sai}.
    \item d) \textbf{Đúng}.
\end{itemize}
}
\end{ex}

\textbf{Phần 3. Trả lời ngắn}
\setcounterpageref{ex}{0}

\begin{ex} %[2D12V62]
Tính hiệu: $813 - 457$.
\shortans{356}
\loigiai{
\textbf{Đáp số:} $356$\\
$813 - 457 = 356$.
}
\end{ex}

\begin{ex} %[2D12V62]
Tìm $x$ biết $904 - x = 368$.
\shortans{536}
\loigiai{
\textbf{Đáp số:} $536$\\
$x = 904 - 368 = 536$.
}
\end{ex}

\begin{ex} %[2D12V62]
Một nông trường có 640 cây ăn quả, trong đó có 275 cây xoài, còn lại là cây nhãn. Có bao nhiêu cây nhãn?
\shortans{365}
\loigiai{
\textbf{Đáp số:} $365$\\
$640 - 275 = 365$ cây.
}
\end{ex}

\begin{ex} %[2D12V62]
Tính giá trị biểu thức: $1 000 - 458 - 242$.
\shortans{300}
\loigiai{
\textbf{Đáp số:} $300$\\
$1 000 - 700 = 300$.
}
\end{ex}

\begin{ex} %[2D12V62]
Số bị trừ là 720, hiệu là 385. Số trừ là bao nhiêu?
\shortans{335}
\loigiai{
\textbf{Đáp số:} $335$\\
$720 - 385 = 335$.
}
\end{ex}

\begin{ex} %[2D12V62]
Một cuộn dây dài 500 m. Người ta cắt một đoạn dài 175 m và một đoạn dài 125 m. Cuộn dây còn lại bao nhiêu mét?
\shortans{200}
\loigiai{
\textbf{Đáp số:} $200$\\
$500 - (175 + 125) = 200$ m.
}
\end{ex}
